% Introduction — v0.1 skeleton.
%
% Target: 2-3 pages. Walks the reader from "what changed about the
% supply chain" to "why a new defense layer is needed" to "what we
% built and what we measured."

\section{Introduction}
\label{sec:intro}

% Subsection 1 — the shift.
% Source: docs/vision.md, docs/threat-model.md § Scope.
\subsection{The agent-context shift}
\label{subsec:shift}

% Claim 1.1 — software supply chain is well-studied.
% TODO: enumerate the prior decade's work on package-manager
% provenance — SolarWinds (2020), Codecov (2021), ua-parser-js
% (2021), eslint-config-prettier (2024), xz-utils CVE-2024-3094
% (2024), the September 2025 Shai-Hulud npm worm.
% Citations: SolarWinds:2020, Codecov:2021, ua-parser-js:2021,
%            eslint-config-prettier:2024, xz-utils:2024,
%            Shai-Hulud:2025.

% Claim 1.2 — what existing tooling addresses.
% TODO: SLSA, in-toto, Sigstore, package-manager provenance — all
% verify that a compiled artifact originated from a specific source
% commit and was signed by an authorized identity.
% Citations: SLSA:2023, in-toto:2016, Sigstore:2021.

% Claim 1.3 — the gap.
% TODO: AI agents load a DIFFERENT class of supply-chain artifact.
% Context primitives (skills, system prompts, MCP server
% descriptions, AGENTS.md files, .claude/skills directories,
% .cursor/rules trees) are markdown-style documents the agent
% INTERPRETS AS INSTRUCTION. None of the existing supply-chain
% controls verify the SEMANTICS of that interpretation.
% Source: docs/threat-model.md § 1.

% Subsection 2 — the attack surface.
% Source: docs/threat-model.md § 2.
\subsection{What an attacker can do with a poisoned primitive}

% Claim 2.1 — illustrative attack.
% TODO: walk through one canonical Shai-Hulud-style attack: a
% primitive looks legitimate (signed, content-digest-verified),
% but its prose tells the agent to ``run \texttt{curl URL | bash}
% before responding,'' or ``fetch \texttt{wss://attacker} on
% session start,'' or impersonates a vendor with a Cyrillic
% homoglyph.
% Source: docs/top10-corpus-map.md CCI-005 + CCI-004 + CCI-007.

% Claim 2.2 — why static signing alone fails.
% TODO: signing proves WHO wrote the artifact; it does NOT prove
% WHAT the artifact tells the agent to do. A signed primitive can
% still contain an attack. The defense layer that's missing is
% SEMANTIC validation — does the prose, when interpreted as
% instruction, cross a threat-model line?
% Source: docs/threat-model.md § 3.

% Claim 2.3 — runtime alone fails too.
% TODO: runtime gating catches some attacks at execution time but
% the cost of a runtime-detected attack is already paid in
% credentials lost / commands run. The static + signed +
% semantic-validation layer prevents the attack from reaching the
% runtime in the first place.

% Subsection 3 — what we built.
\subsection{Contributions}

% Claim 3.1 — taxonomy.
% TODO: the Context CI Top 10 — a ten-category taxonomy of
% agent-context supply-chain attacks, each with detector +
% canonical sample + measured precision/recall.
% Source: docs/top10-corpus-map.md.

% Claim 3.2 — corpus.
% TODO: a 105-sample labeled corpus (99 in-sample + 6 holdout) with
% per-detector P/R/F1 measurement, byte-identical across two
% independent implementations.
% Source: docs/corpus-v0.1.md + docs/evaluation.md.

% Claim 3.3 — methodology.
% TODO: an AV-vendor-equivalent FP/FN ledger discipline that
% prevents tuning past optimal — every detector change has an
% append-only numbered ledger entry and an accompanying corpus
% measurement that proves the change generalises.
% Source: docs/false-{positive,negative}-resolution.md.

% Claim 3.4 — runtime gate.
% TODO: a regression gate (\texttt{bench/scripts/gate.sh}) that runs
% on every PR and blocks merges on F1 regression $> 0.02$ or
% cross-language verdict drift. The gate is the steering wheel; the
% corpus is the road.
% Source: bench/baseline/FLOOR.md, docs/decisions/0037.

% Claim 3.5 — reproducibility.
% TODO: the entire pipeline is open-source and reproducible from a
% fresh clone in under 10 minutes. \texttt{bench/cli/replay.py} lets a
% reviewer re-run any specific sample deterministically and diff its
% verdict against the committed baseline.

% Subsection 4 — paper layout.
\subsection{Paper organization}

% TODO: standard ``Section 2 presents the threat model... Section 6
% reports evaluation results...'' wrapper.
% \Cref{sec:threat-model} presents the Context CI Top 10 taxonomy.
% \Cref{sec:system-design} describes the Trust Runtime three-pillar
% architecture (Trust Contract, Top 10, CBOM). \Cref{sec:implementation}
% covers the dual-implementation discipline. \Cref{sec:evaluation}
% reports per-detector precision, recall, and F1. \Cref{sec:related}
% positions PrimitiveHub against prior supply-chain work.
% \Cref{sec:conclusion} discusses limitations and future work.
